Every function from the five-element field to itself is represented by exactly one polynomial of degree at most $4$: existence follows by interpolation, and uniqueness follows because a nonzero polynomial of degree at most $4$ cannot have five roots. It suffices to count functions $f$ with $f^3=f$.

Let $I$ be the image, of size $k$. For $y=f(x)\in I$, the condition says $f^2(y)=y$. Thus the restriction to $I$ is an involution, and every point outside $I$ may independently map to any point of $I$. Conversely all such choices work and have image exactly $I$.

The involution counts on $k=1,\ldots,5$ points are $1,2,4,10,26$, obtained from $J_k=J_{k-1}+(k-1)J_{k-2}$ with $J_0=J_1=1$ (the first point is fixed or paired). The desired total is
\[
\sum_{k=1}^5\binom5kJ_k k^{5-k}
=5+160+360+200+26=\boxed{751}.
\]
